%%%%%%%%%%%%%%%%%%%%%%%%%%%%%%%%%%%%%%%%%%%%%%%%%%%%%%%%%%%%%%%%%%%%%%%%%%%%%%%%
%2345678901234567890123456789012345678901234567890123456789012345678901234567890
%        1         2         3         4         5         6         7         8

\documentclass[letterpaper, 10 pt, conference]{ieeeconf}  % Comment this line out if you need a4paper

%\documentclass[a4paper, 10pt, conference]{ieeeconf}      % Use this line for a4 paper

\IEEEoverridecommandlockouts                              % This command is only needed if
                                                          % you want to use the \thanks command

\overrideIEEEmargins                                      % Needed to meet printer requirements.


\input{usepackage}
% Extra packages required for this milestone (table, figures, hyperlinks).
% Loaded explicitly so this file compiles even if usepackage.sty does not
% already pull them in. graphicx is harmless if duplicated.
\usepackage{graphicx}
\usepackage{booktabs}
\usepackage[caption=false,font=footnotesize]{subfig}
\usepackage{hyperref}
\usepackage{xcolor}
\hypersetup{colorlinks=true, urlcolor=blue, linkcolor=black, citecolor=black}

% Tighten floats so the 1-page limit holds with table + figure block.
\setlength{\textfloatsep}{4pt}
\setlength{\intextsep}{4pt}
\setlength{\abovecaptionskip}{2pt}
\setlength{\belowcaptionskip}{0pt}

% Centralise the figure search path so the date-stamped folder name does not
% trip chktex warning 8 on every \includegraphics line. On Overleaf, upload
% the six PNGs into submission/week6-7/figs/ and change this to {{figs/}}.
\graphicspath{{../../sim_output/20260531-192010_hard_brake_GOOD/}} % chktex 8

\title{\LARGE \bf
Week 6-7 Second Attempt: CVaR-MPPI for Highway Lane-Change Planning}


\author{\parbox{4 in}{\centering Ayush Agrawal\\ UID: 306184708}
\thanks{UCLA ECE M237 Spring 2026 Course Project Milestone Reports. Week 6-7 Second Attempt.}
}

\begin{document}



\maketitle

\section{Summary of your initial idea and results}
I augment Model Predictive Path Integral (MPPI)~\cite{mppi} with a Conditional Value-at-Risk (CVaR)~\cite{rockafellar} objective. The goal is to suppress maneuvers that are efficient on average but unsafe under rare traffic outcomes. The constrained problem is $\min_{v}\mathbb{E}[\phi(x_K)+\sum_{k}(q(x_k)+\tfrac{\lambda}{2}v_k^{\!\top}\Sigma_\epsilon^{-1}v_k)]$ s.t.\ $\mathrm{CVaR}_\alpha(L(\tilde x))\!\le\!C_u$. I decompose it into two stages. Offline, a CVaR-MDP value iteration~\cite{chow_cvarmdp} supplies a tail-risk-shaped terminal cost $V^{*}(x,y)$. Online, a Risk-Aware MPPI~\cite{ramppi} enforces the tail constraint via a per-rollout Monte-Carlo CVaR penalty. Week-4 results reported a tuned \emph{risk-neutral} MPPI baseline in HighwayEnv~\cite{highwayenv}. The baseline avoided crashes on benign traffic. It still violated comfort and safety limits with large $|a_y|$ and low $\mathrm{TTC}$. These violations motivate the CVaR objective.

\section{Summary of feedback from peer reviews and grade}
The peer reviewer raised five concerns. (i) State the novelty more sharply relative to RA-MPPI~\cite{ramppi}. (ii) Tighten scope by prioritising the online CVaR-MPPI; defer the offline $V^{*}$ as a stretch goal. (iii) Define $L(\tilde x)$, the empirical $\widehat{\mathrm{VaR}}/\widehat{\mathrm{CVaR}}$ estimator, the number $N$ of disturbance rollouts, and the penalty weight $A$ explicitly. (iv) Clarify whether the traffic predictor is interactive or exogenous. (v) Add a clean ablation between risk-neutral MPPI, online CVaR-MPPI, and CVaR-MPPI with $V^{*}$ terminal. Example-side comments asked for multi-agent reactive scenarios and a richer multi-seed metric set (collision rate, near-miss rate, $\mathrm{TTC}_{\min}$, $|a_y|$, $|j_y|$, $T_{LC}$, success rate). A notation inconsistency ($M\!=\!128$ vs $M\!=\!256$) was also flagged. The Gradescope grade reinforced the central message. Stress-test the controller. Design experiments that prove the value of the CVaR penalty.

\section{Updated method}
I adopted the reviewer's focused scope and implemented \textbf{Track~A}: the online CVaR-MPPI of~\cite{ramppi}. Offline $V^{*}$ is deferred. The controller subclasses the risk-neutral MPPI sampler. It modifies only the cost aggregator. For each of $M\!=\!256$ sampled control sequences $u^{m}$, the planner draws $N\!=\!24$ disturbed traffic futures via $\tilde F$. $\tilde F$ is the constant-velocity nominal plus zero-mean Gaussian acceleration noise $\sigma_a\!=\!1.5\,\mathrm{m/s^2}$, so $\mathbb{E}_w[\tilde F]\!=\!F$. I then evaluate the trajectory risk $L^{m,n}\!=\!\sum_{k}\ell(\tilde x^{m,n}_k)$. The Lagrangian penalty $J_C^{m}\!=\!A\,\widehat{\mathrm{CVaR}}_\alpha(L^{m,*})\,\mathbf{1}\{\widehat{\mathrm{CVaR}}_\alpha\!>\!C_u\}$ is added before the softmin, with defaults $\alpha\!=\!0.9$, $C_u\!=\!4.0$, $A\!=\!80$. $\widehat{\mathrm{CVaR}}_\alpha$ is the tail-average estimator over the worst $\lceil N(1\!-\!\alpha)\rceil$ samples~\cite{rockafellar}. The risk cost $\ell(\tilde x)$ is an explicit TTC-shortfall potential plus a lateral-encroachment indicator. It is distinct from the task cost $q$. The traffic predictor $\tilde F$ is exogenous (independent of $u^{m}$), matching~\cite{ramppi}. The $M\!=\!128/256$ inconsistency is resolved (unified $M\!=\!256$).

\section{Updated working example}
I replaced the single-episode Week-4 test with a paired-seed head-to-head between risk-neutral MPPI and CVaR-MPPI\@. The harness covers three scenarios. A hard-brake leader at $-10\,\mathrm{m/s^2}$ with $8\,\mathrm{m}$ headway. An aggressive cut-in. A natural lane change. Each scenario runs over $\ge\!20$ seeds. I report collision rate, near-miss rate ($\mathrm{TTC}\!<\!1.5\,\mathrm{s}$), $\mathrm{TTC}_{\min}$, $|a_y|_{\mathrm{peak}}$, $|j_y|_{\mathrm{peak}}$, $T_{LC}$, and $e_{y,\mathrm{RMS}}$. The hard-brake case is a two-lane shift (lane~$2\!\to\!$lane~$0$) executed simultaneously with the leader's emergency brake. This setup is designed expressly to surface the tail-risk failures the reviewer and Gradescope feedback asked me to stress-test.

\section{Results}
After tuning (Table~\ref{tab:updates}), the risk-neutral MPPI baseline \emph{completes the harder stress test without collision} ($160$ steps, $7.95\,\mathrm{s}$ simulated, seed 3; video: \href{https://drive.google.com/REPLACE-WITH-YOUR-LINK}{run.gif}). The achieved maneuver still exceeds the comfort and safety thresholds. Specifically, $|a_y|_{\mathrm{peak}}\!=\!23.3\,\mathrm{m/s^2}$, $|j_y|_{\mathrm{peak}}\!=\!97.5\,\mathrm{m/s^3}$, $\mathrm{TTC}_{\min}\!=\!0.50\,\mathrm{s}$, and $T_{LC}\!=\!6.45\,\mathrm{s}$. This is exactly the gap CVaR-MPPI is designed to close. The new outcome differs from Week 4 in two ways. First, the baseline now survives a strictly harder scenario. Any improvement attributable to CVaR will therefore not be masked by baseline failure. Second, the full Track-A infrastructure (\texttt{risk.py}, \texttt{cvar.py}, \texttt{stochastic\_idm.py}, \texttt{cvar\_mppi.py}, \texttt{compare\_controllers.py}) is implemented and unit-tested. The paired CVaR-vs-MPPI ablation is queued.

\begin{table}[h]
\centering\footnotesize
\caption{MPPI updates yielding the no-crash hard-brake maneuver.}
\label{tab:updates}
\begin{tabular}{@{}lccl@{}}
\toprule
\textbf{Knob} & \textbf{Week 4} & \textbf{Updated} & \textbf{Effect} \\
\midrule
Temperature $\lambda$            & 1.0   & 2.0   & sharper softmin \\
Accel.\ noise $\sigma_a$         & 1.5   & 0.6   & smoother controls \\
Steer noise $\sigma_\delta$      & 0.08  & 0.03  & damps oscillation \\
Heading cost $w_\psi$            & none  & 18.0  & lateral-oscillation damping \\
Time-headway $w_{\mathrm{thw}}$  & none  & 10.0  & enforces 2-s rule \\
Road barrier $w_{\mathrm{road}}$ & none  & 500.0 & prevents off-road excursion \\
Bicycle integrator               & Euler ($L$) & slip-$\beta$ ($L/2$) & matches HighwayEnv \\
Samples $M$                      & 128/256 & 256 & resolves notation \\
\bottomrule
\end{tabular}
\end{table}

\begin{figure*}[!htbp]
\centering
\subfloat[Top-down $X$-$Y$ trajectory with traffic snapshots. The ego completes the two-lane shift (lane $2\to$lane $0$) while the leader executes a $-10\,\mathrm{m/s^2}$ brake. The path remains crash free.]{%
  \includegraphics[width=0.32\textwidth]{trajectory_overview.png}\label{fig:traj}}\hfill
\subfloat[Lateral profile $Y(t)$ and tracking error $e_y(t)$ vs the target lane. The settled error $e_{y,\mathrm{RMS}}\!=\!0.110\,\mathrm{m}$ sits just above the $0.10\,\mathrm{m}$ acceptable band.]{%
  \includegraphics[width=0.32\textwidth]{lateral_profile.png}\label{fig:lat}}\hfill
\subfloat[$|a_y|(t)$ and $|j_y|(t)$ against ISO 11270 comfort bands. Both peaks exceed the hard limit, marking the comfort gap the CVaR penalty is designed to close.]{%
  \includegraphics[width=0.32\textwidth]{comfort.png}\label{fig:comfort}}

\vspace{3pt}
\subfloat[$\mathrm{TTC}(t)$ against ISO 17387 bands. The minimum TTC of $0.50\,\mathrm{s}$ sits below the $2\,\mathrm{s}$ warning threshold and pinpoints the tail event CVaR-MPPI must suppress.]{%
  \includegraphics[width=0.32\textwidth]{safety.png}\label{fig:safety}}\hfill
\subfloat[Commanded acceleration $a$ and steering $\delta$ versus time. The tightened noise variances $(\sigma_a,\sigma_\delta)\!=\!(0.6, 0.03)$ keep the actuation bounded throughout the brake event.]{%
  \includegraphics[width=0.32\textwidth]{controls.png}\label{fig:ctrl}}\hfill
\subfloat[Rollout cost histogram with $\widehat{\mathrm{VaR}}_\alpha$ and $\widehat{\mathrm{CVaR}}_\alpha$ markers. The right tail is precisely the mass that the Lagrangian penalty $J_C^{m}$ will down-weight.]{%
  \includegraphics[width=0.32\textwidth]{cost_distribution.png}\label{fig:cost}}
\caption{Diagnostics for the no-crash hard-brake maneuver of Sec.\,V (seed 3, $160$ steps, $7.95\,\mathrm{s}$). Six panels track trajectory, lateral tracking, ISO comfort, safety (TTC), commanded controls, and the rollout cost distribution.}\label{fig:hard_brake_diagnostics}
\end{figure*}

\section{Remaining Steps}
Three steps remain. (i) Run the 20-seed paired head-to-head (\texttt{compare\_controllers.py}) on \texttt{hard\_brake}, \texttt{cutin}, and \texttt{lane\_change}; tabulate collision and near-miss rates with comfort metrics. (ii) Sweep $(\alpha, C_u, A)$ via \texttt{cvar\_sweep.py} for the safety vs.\ performance frontier. (iii) If time permits, prototype an offline $V^{*}$ on a tiny $|\mathcal{X}|\!\sim\!10^{3}$, $|\mathcal{Y}|\!=\!21$ grid as a third arm of the ablation.

% --------------------------------------------------------------------------
% Bibliography. Two options:
%   (A) inline thebibliography (used below) -- self-contained, no .bib needed.
%   (B) external .bib via \bibliography{IEEEabrv,references} with
%       \bibliographystyle{IEEEtran} -- matches the formal template default.
% Switch to (B) once a references.bib file exists in this folder.
% --------------------------------------------------------------------------
\begin{thebibliography}{9}\footnotesize
\bibitem{mppi} G.~Williams, A.~Aldrich, E.~A.~Theodorou, ``Model predictive path integral control: From theory to parallel computation,'' \emph{J.~Guid.\ Control Dyn.}, 40(2):344-357, 2017.
\bibitem{rockafellar} R.~T.~Rockafellar, S.~Uryasev, ``Optimization of conditional value-at-risk,'' \emph{J.~Risk}, 2(3):21-41, 2000.
\bibitem{chow_cvarmdp} Y.~Chow, A.~Tamar, S.~Mannor, M.~Pavone, ``Risk-sensitive and robust decision-making: a CVaR optimization approach,'' in \emph{NeurIPS} 28, 2015.
\bibitem{ramppi} J.~Yin, Z.~Zhang, P.~Tsiotras, ``Risk-aware MPPI control using conditional value-at-risk,'' arXiv:2209.12842, 2022.
\bibitem{highwayenv} E.~Leurent, ``An environment for autonomous driving decision-making,'' \texttt{github.com/Farama-Foundation/HighwayEnv}, 2018.
\end{thebibliography}



\end{document}
